\FigureLayoutDeclare{img/2019/national_paper_1_science/q6_fig1.png}{4.0cm}{0bp 0bp 0bp 0bp}

\chapter{2019年全国I卷（理）}

\section{单选题}

\begin{problem}
已知集合 \(M=\{x\mid -4<x<2\}\)，\(N=\{x\mid x^2-x-6<0\}\)，则 \(M\cap N=\)（\quad）

\choices
    {\(\{x\mid -4<x<3\}\)}
    {\(\{x\mid -4<x<-2\}\)}
    {\(\{x\mid -2<x<2\}\)}
    {\(\{x\mid 2<x<3\}\)}

\begin{answer}
C
\end{answer}
\end{problem}

\begin{problem}
设复数 \(z\) 满足 \(|z-\i|=1\)，\(z\) 在复平面内对应的点为 \((x,y)\)，则（\quad）

\choices
    {\((x+1)^2+y^2=1\)}
    {\((x-1)^2+y^2=1\)}
    {\(x^2+(y-1)^2=1\)}
    {\(x^2+(y+1)^2=1\)}

\begin{answer}
C
\end{answer}
\end{problem}

\begin{problem}
已知 \(a=\log_{2}0.2\)，\(b=2^{0.2}\)，\(c=0.2^{0.3}\)，则（\quad）

\choices
    {\(a<b<c\)}
    {\(a<c<b\)}
    {\(c<a<b\)}
    {\(b<c<a\)}

\begin{answer}
B
\end{answer}
\end{problem}

\begin{problem}
古希腊时期，人们认为最美人体的头顶至肚脐的长度与肚脐至足底的长度之比是 \(\dfrac{\sqrt5-1}{2}\approx0.618\)（称为黄金分割比例）. 著名的“断臂维纳斯”便是如此. 此外，最美人体的头顶至咽喉的长度与咽喉至肚脐的长度之比也是 \(\dfrac{\sqrt5-1}{2}\). 若某人满足上述两个黄金分割比例，且腿长为 \(105\text{ cm}\)，头顶至脖子下端的长度为 \(26\text{ cm}\)，则其身高可能是（\quad）

\bitmapfigure[max width=5.0cm]{img/2019/national_paper_1_science/q4_fig1.png}

\choices
    {\(165\text{ cm}\)}
    {\(175\text{ cm}\)}
    {\(185\text{ cm}\)}
    {\(190\text{ cm}\)}

\begin{answer}
B
\end{answer}
\end{problem}

\begin{problem}
函数 \(f(x)=\dfrac{\sin x+x}{\cos x+x^2}\) 在 \([-\pi,\pi]\) 的图象大致为（\quad）

\choices
    {\(\choicebitmap[width=0.15\paperwidth]{img/2019/national_paper_1_science/q5_choice_a.png}\)}
    {\(\choicebitmap[width=0.15\paperwidth]{img/2019/national_paper_1_science/q5_choice_b.png}\)}
    {\(\choicebitmap[width=0.15\paperwidth]{img/2019/national_paper_1_science/q5_choice_c.png}\)}
    {\(\choicebitmap[width=0.15\paperwidth]{img/2019/national_paper_1_science/q5_choice_d.png}\)}

\begin{answer}
D
\end{answer}
\end{problem}

\begin{problem}
我国古代典籍《周易》用“卦”描述万物的变化. 每一“重卦”由从下到上排列的 \(6\) 个爻组成，爻分为阳爻“——”和阴爻“— —”. 如图就是一“重卦”. 在所有重卦中随机取一重卦，则该重卦恰有 \(3\) 个阳爻的概率是（\quad）

\bitmapfigure[max width=6.0cm]{img/2019/national_paper_1_science/q6_fig1.png}

\choices
    {\(\dfrac{5}{16}\)}
    {\(\dfrac{3}{8}\)}
    {\(\dfrac{1}{2}\)}
    {\(\dfrac{1}{16}\)}

\begin{answer}
A
\end{answer}
\end{problem}

\begin{problem}
已知非零向量 \(\bs{a}\)，\(\bs{b}\) 满足 \(|\bs{a}|=2|\bs{b}|\)，且 \((\bs{a}-\bs{b})\perp \bs{b}\)，则 \(\bs{a}\) 与 \(\bs{b}\) 的夹角为（\quad）

\choices
    {\(\dfrac{\pi}{6}\)}
    {\(\dfrac{\pi}{3}\)}
    {\(\dfrac{2\pi}{3}\)}
    {\(\dfrac{5\pi}{6}\)}

\begin{answer}
B
\end{answer}
\end{problem}

\begin{problem}
如图是求 \(\dfrac{1}{2+\dfrac{1}{2+\dfrac{1}{2}}}\) 的程序框图，图中空白框中应填入（\quad）

\bitmapfigure[max width=6.0cm]{img/2019/national_paper_1_science/q8_fig1.png}

\choices
    {\(A=\dfrac{1}{2+A}\)}
    {\(A=2+\dfrac{1}{A}\)}
    {\(A=\dfrac{1}{1+2A}\)}
    {\(A=1+\dfrac{1}{2A}\)}

\begin{answer}
A
\end{answer}
\end{problem}

\begin{problem}
设 \(S_n\) 为等差数列 \(\{a_n\}\) 的前 \(n\) 项和. 已知 \(S_4=0\)，\(a_5=5\)，则（\quad）

\choices
    {\(a_n=2n-5\)}
    {\(a_n=3n-10\)}
    {\(S_n=2n^2-8n\)}
    {\(S_n=n^2-2n\)}

\begin{answer}
A
\end{answer}
\end{problem}

\begin{problem}
已知椭圆 \(C\) 的焦点为 \(F_1(-1,0)\)，\(F_2(1,0)\)，过 \(F_2\) 的直线与 \(C\) 交于 \(A\)，\(B\) 两点. 若 \(|AF_2|=2|F_2B|\)，\(|AB|=|BF_1|\)，则 \(C\) 的方程为（\quad）

\choices
    {\(\dfrac{x^2}{2}+y^2=1\)}
    {\(\dfrac{x^2}{3}+\dfrac{y^2}{2}=1\)}
    {\(\dfrac{x^2}{4}+\dfrac{y^2}{3}=1\)}
    {\(\dfrac{x^2}{5}+\dfrac{y^2}{4}=1\)}

\begin{answer}
B
\end{answer}
\end{problem}

\begin{problem}
关于函数 \(f(x)=\sin|x|+|\sin x|\) 有下述四个结论：

\begin{circlelist}
\item \(f(x)\) 是偶函数；

\item \(f(x)\) 在区间 \(\left(\dfrac{\pi}{2},\pi\right)\) 单调递增；

\item \(f(x)\) 在 \([-\pi,\pi]\) 有 \(4\) 个零点；

\item \(f(x)\) 的最大值为 \(2\).
\end{circlelist}

其中所有正确结论的编号是（\quad）

\choices
    {\(\circled{1}\circled{2}\circled{4}\)}
    {\(\circled{2}\circled{4}\)}
    {\(\circled{1}\circled{4}\)}
    {\(\circled{1}\circled{3}\)}

\begin{answer}
C
\end{answer}
\end{problem}

\begin{problem}
已知三棱锥 \(P-ABC\) 的四个顶点在球 \(O\) 的球面上，\(PA=PB=PC\)，\(\triangle ABC\) 是边长为 \(2\) 的正三角形，\(E\)，\(F\) 分别是 \(PA\)，\(AB\) 的中点，\(\angle CEF=90^\circ\)，则球 \(O\) 的体积为（\quad）

\choices
    {\(8\sqrt6\pi\)}
    {\(4\sqrt6\pi\)}
    {\(2\sqrt6\pi\)}
    {\(\sqrt6\pi\)}

\begin{answer}
D
\end{answer}
\end{problem}

\section{填空题}

\begin{problem}
曲线 \(y=3(x^2+x)\e^x\) 在点 \((0,0)\) 处的切线方程为\fillinblank{}.

\begin{answer}
\(y=3x\)
\end{answer}
\end{problem}

\begin{problem}
记 \(S_n\) 为等比数列 \(\{a_n\}\) 的前 \(n\) 项和. 若 \(a_1=\dfrac13\)，\(a_4^2=a_6\)，则 \(S_5=\fillinblank{}\).

\begin{answer}
\(\dfrac{121}{3}\)
\end{answer}
\end{problem}

\begin{problem}
甲、乙两队进行篮球决赛，采取七场四胜制（当一队赢得四场胜利时，该队获胜，决赛结束）. 根据前期比赛成绩，甲队的主客场安排依次为“主主客客主客主”. 设甲队主场取胜的概率为 \(0.6\)，客场取胜的概率为 \(0.5\)，且各场比赛结果相互独立，则甲队以 \(4:1\) 获胜的概率是\fillinblank{}.

\begin{answer}
\(0.18\)
\end{answer}
\end{problem}

\begin{problem}
已知双曲线 \(C:\dfrac{x^2}{a^2}-\dfrac{y^2}{b^2}=1\)（\(a>0\)，\(b>0\)）的左、右焦点分别为 \(F_1\)，\(F_2\). 过 \(F_1\) 的直线与 \(C\) 的两条渐近线分别交于 \(A\)，\(B\) 两点，若 \(\overrightarrow{F_1A}=\overrightarrow{AB}\)，\(\overrightarrow{F_1B}\cdot\overrightarrow{F_2B}=0\)，则 \(C\) 的离心率为\fillinblank{}.

\begin{answer}
2
\end{answer}
\end{problem}

\section{解答题}

\begin{problem}
\(\triangle ABC\) 的内角 \(A\)，\(B\)，\(C\) 的对边分别为 \(a\)，\(b\)，\(c\). 设 \((\sin B-\sin C)^2=\sin^2A-\sin B\sin C\).

\begin{enumerate}
\item 求 \(A\)；

\item 若 \(\sqrt2\,a+b=2c\)，求 \(\sin C\).
\end{enumerate}
\end{problem}

\begin{solution}
\(A=60^{\circ}\)；\(\sin C=\dfrac{\sqrt6+\sqrt2}{4}\).
\end{solution}

\begin{problem}
如图，直四棱柱 \(ABCD-A_1B_1C_1D_1\) 的底面是菱形，\(AA_1=4\)，\(AB=2\)，\(\angle BAD=60^\circ\). \(E\)，\(M\)，\(N\) 分别是 \(BC\)，\(BB_1\)，\(A_1D\) 的中点.

\begin{enumerate}
\item 证明：\(MN\parallel\) 平面 \(C_1DE\)；

\item 求二面角 \(A-MA_1-N\) 的正弦值.

\end{enumerate}
\bitmapfigure[max width=5.1cm]{img/2019/national_paper_1_science/q18_fig1.png}
\end{problem}

\begin{solution}
\(MN\parallel\) 平面 \(C_1DE\)；二面角 \(A-MA_1-N\) 的正弦值为 \(\dfrac{\sqrt{10}}{5}\).
\end{solution}

\begin{problem}
已知抛物线 \(C:y^2=3x\) 的焦点为 \(F\)，斜率为 \(\dfrac32\) 的直线 \(l\) 与 \(C\) 的交点为 \(A\)，\(B\)，与 \(x\) 轴的交点为 \(P\).

\begin{enumerate}
\item 若 \(|AF|+|BF|=4\)，求 \(l\) 的方程；

\item 若 \(\overrightarrow{AP}=3\overrightarrow{PB}\)，求 \(|AB|\).
\end{enumerate}
\end{problem}

\begin{solution}
\(l\) 的方程为 \(y=\dfrac{3}{2}x-\dfrac{7}{8}\)；\(|AB|=\dfrac{4\sqrt{13}}{3}\).
\end{solution}

\begin{problem}
已知函数 \(f(x)=\sin x-\ln(1+x)\)，且 \(f'(x)\) 为 \(f(x)\) 的导数. 证明：

\begin{enumerate}
\item \(f'(x)\) 在区间 \(\left(-1,\dfrac{\pi}{2}\right)\) 存在唯一极大值点；

\item \(f(x)\) 有且仅有 \(2\) 个零点.
\end{enumerate}
\end{problem}

\begin{solution}
\(f'(x)\) 在区间 \(\left(-1,\dfrac{\pi}{2}\right)\) 存在唯一极大值点；\(f(x)\) 有且仅有 \(2\) 个零点.
\end{solution}

\begin{problem}
为治疗某种疾病，研制了甲、乙两种新药，希望知道哪种新药更有效，为此进行动物试验. 试验方案如下：每一轮选取两只白鼠对药效进行对比试验. 对于两只白鼠，随机选一只施以甲药，另一只施以乙药. 一轮的治疗结果得出后，再安排下一轮试验. 当其中一种药治愈的白鼠比另一种药治愈的白鼠多 \(4\) 只时，就停止试验，并认为治愈只数多的药更有效. 为了方便描述问题，约定：对于每轮试验，若施以甲药的白鼠治愈且施以乙药的白鼠未治愈则甲药得 \(1\) 分，乙药得 \(-1\) 分；若施以乙药的白鼠治愈且施以甲药的白鼠未治愈则乙药得 \(1\) 分，甲药得 \(-1\) 分；若都治愈或都未治愈则两种药均得 \(0\) 分. 甲、乙两种药的治愈率分别记为 \(\alpha\) 和 \(\beta\)，一轮试验中甲药的得分记为 \(X\).

\begin{enumerate}
\item 求 \(X\) 的分布列；

\item 若甲药、乙药在试验开始时都赋予 \(4\) 分，\(p_i\)（\(i=0\)，\(1\)，\(\cdots\)，\(8\)）表示“甲药的累计得分为 \(i\) 时，最终认为甲药比乙药更有效”的概率，则 \(p_0=0\)，\(p_8=1\)，\(p_i=ap_{i-1}+bp_i+cp_{i+1}\)（\(i=1\)，\(2\)，\(\cdots\)，\(7\)），其中 \(a=P(X=-1)\)，\(b=P(X=0)\)，\(c=P(X=1)\). 假设 \(\alpha=0.5\)，\(\beta=0.8\).

\begin{enumerate}
\item 证明：\(\{p_{i+1}-p_i\}\)（\(i=0\)，\(1\)，\(2\)，\(\cdots\)，\(7\)）为等比数列；
\item 求 \(p_4\)，并根据 \(p_4\) 的值解释这种试验方案的合理性.
\end{enumerate}
\end{enumerate}
\end{problem}

\begin{solution}
\(P(X=-1)=(1-\alpha)\beta\)，\(P(X=0)=\alpha\beta+(1-\alpha)(1-\beta)\)，\(P(X=1)=\alpha(1-\beta)\)；\(p_4=\dfrac{1}{257}\).
\end{solution}

\section{选考题：解答题}

\begin{problem}
在直角坐标系 \(xOy\) 中，曲线 \(C\) 的参数方程为 \(\left\{\begin{aligned}x&=\dfrac{1-t^2}{1+t^2},\\y&=\dfrac{4t}{1+t^2},\end{aligned}\right.\)（\(t\) 为参数）. 以原点为极点，\(x\) 轴正半轴为极轴建立极坐标系，直线 \(l\) 的极坐标方程为 \(2\rho\cos\theta+\sqrt3\,\rho\sin\theta+11=0\).

\begin{enumerate}
\item 求 \(C\) 的直角坐标方程与 \(l\) 的直角坐标方程；

\item 求 \(C\) 上的点到 \(l\) 的最小距离.
\end{enumerate}
\end{problem}

\begin{solution}
\(C\) 的直角坐标方程为 \(x^2+\dfrac{y^2}{4}=1\)，\(l\) 的直角坐标方程为 \(2x+\sqrt3y+11=0\)；\(C\) 上的点到 \(l\) 的最小距离为 \(\sqrt7\).
\end{solution}

\begin{problem}
已知 \(a\)，\(b\)，\(c\) 为正数，且满足 \(abc=1\). 证明：

\begin{enumerate}
\item \(\dfrac{1}{a}+\dfrac{1}{b}+\dfrac{1}{c}\le a^2+b^2+c^2\)；

\item \((a+b)^3+(b+c)^3+(c+a)^3\ge24\).
\end{enumerate}
\end{problem}

\begin{solution}
两不等式均成立.
\end{solution}
